% Môn: Vật lý
% Lớp: 12
% Chương: Chương II. Khí lí tưởng
% Bài: L12C2 Bài 10. Định luật Charles · Tính áp suất, thể tích trong quá trình đẳng nhiệt
% ===== Câu 2 =====
% ID: L12C2B10-144-TL
% Mức: VD
\dangbt{Tính áp suất, thể tích trong quá trình đẳng nhiệt}
\begin{bt}
% ID: L12C2B10-144-TL
% Mức: VD
Trong dạng “Bài toán tổng hợp về quá trình đẳng nhiệt”, Thể tích giảm từ 28 lít còn 7 lít ở nhiệt độ không đổi. Áp suất tăng bao nhiêu lần?
\loigiai{
$p_2/p_1=V_1/V_2=4$.
}
\end{bt}

% ===== Câu 3 =====
% ID: L12C2B10-145-TLN
% Mức: TH
\begin{ex}
% ID: L12C2B10-145-TLN
% Mức: TH
Trong dạng “Bài toán tổng hợp về quá trình đẳng nhiệt”, Khí 6 lít ở $100\,\text{kPa}$ được nén đẳng nhiệt còn 3 lít. Tính áp suất sau theo kPa.
\shortans{200}
\loigiai{
$p_2=100\cdot6/3=200$ kPa.
}
\end{ex}

% ===== Câu 6 =====
% ID: L12C2B10-146-TL
% Mức: VD
\begin{bt}
% ID: L12C2B10-146-TL
% Mức: VD
Trong dạng “Bài toán tổng hợp về quá trình đẳng nhiệt”, 130
\loigiai{
$p=390/3=130$ kPa.
}
\end{bt}

% ===== Câu 7 =====
% ID: L12C2B10-147-TLN
% Mức: TH
\begin{ex}
% ID: L12C2B10-147-TLN
% Mức: TH
Trong dạng “Bài toán tổng hợp về quá trình đẳng nhiệt”, Khí 5 lít ở $330\,\text{kPa}$ giãn đẳng nhiệt đến áp suất $110\,\text{kPa}$. Tính thể tích sau theo lít.
\shortans{15}
\loigiai{
$V_2=330\cdot5/110=15$ lít.
}
\end{ex}

% ===== Câu 10 =====
% ID: L12C2B10-148-TL
% Mức: VD
\begin{bt}
% ID: L12C2B10-148-TL
% Mức: VD
Trong dạng “Bài toán tổng hợp về quá trình đẳng nhiệt”, $\rho_2/\rho_1=p_2/p_1=3$.
\loigiai{
$\rho_2/\rho_1=p_2/p_1=3$.
}
\end{bt}

% ===== Câu 11 =====
% ID: L12C2B10-149-TLN
% Mức: VD
\begin{ex}
% ID: L12C2B10-149-TLN
% Mức: VD
Trong dạng “Bài toán tổng hợp về quá trình đẳng nhiệt”, Thể tích giảm từ 28 lít còn 7 lít ở nhiệt độ không đổi. Áp suất tăng bao nhiêu lần?
\shortans{4}
\loigiai{
$p_2/p_1=V_1/V_2=4$.
}
\end{ex}

% ===== Câu 12 =====
% ID: L12C2B10-150-TN
% Mức: TH
\begin{ex}
% ID: L12C2B10-150-TN
% Mức: TH
Trong dạng “Bài toán tổng hợp về quá trình đẳng nhiệt”, Khí có áp suất $130\,\text{kPa}$, thể tích 10 lít được nén đẳng nhiệt còn 5 lít. Áp suất sau là
\choice
{\True $260\,\text{kPa}$}
{$65\,\text{kPa}$}
{$130\,\text{kPa}$}
{$520\,\text{kPa}$}
\loigiai{
Đáp án A. $p_2=p_1V_1/V_2=130\cdot10/5=260$ kPa.
}
\end{ex}

% ===== Câu 15 =====
% ID: L12C2B10-151-TLN
% Mức: VD
\begin{ex}
% ID: L12C2B10-151-TLN
% Mức: VD
Trong dạng “Bài toán tổng hợp về quá trình đẳng nhiệt”, Một lượng khí có $pV=390$ kPa.lít. Khi V=3 lít, tính p theo kPa.
\shortans{130}
\loigiai{
$p=390/3=130$ kPa.
}
\end{ex}

% ===== Câu 16 =====
% ID: L12C2B10-152-TN
% Mức: TH
\begin{ex}
% ID: L12C2B10-152-TN
% Mức: TH
Trong dạng “Bài toán tổng hợp về quá trình đẳng nhiệt”, Khí có áp suất $90\,\text{kPa}$, thể tích 7 lít giãn đẳng nhiệt đến 21 lít. Áp suất sau là
\choice
{\True $30\,\text{kPa}$}
{$270\,\text{kPa}$}
{$90\,\text{kPa}$}
{$10\,\text{kPa}$}
\loigiai{
Đáp án A. $p_2=90\cdot7/21=30$ kPa.
}
\end{ex}

% ===== Câu 19 =====
% ID: L12C2B10-153-TLN
% Mức: VD
\begin{ex}
% ID: L12C2B10-153-TLN
% Mức: VD
Trong dạng “Bài toán tổng hợp về quá trình đẳng nhiệt”, Áp suất đẳng nhiệt tăng từ $140\,\text{kPa}$ lên $420\,\text{kPa}$. Khối lượng riêng tăng bao nhiêu lần?
\shortans{3}
\loigiai{
$\rho_2/\rho_1=p_2/p_1=3$.
}
\end{ex}

% ===== Câu 20 =====
% ID: L12C2B10-154-TN
% Mức: TH
\begin{ex}
% ID: L12C2B10-154-TN
% Mức: TH
Trong dạng “Bài toán tổng hợp về quá trình đẳng nhiệt”, Thể tích khí giảm 4 lần trong quá trình đẳng nhiệt thì áp suất
\choice
{\True tăng 4 lần}
{giảm 4 lần}
{không đổi}
{tăng 16 lần}
\loigiai{
Đáp án A. Do $p$ tỉ lệ nghịch với $V$.
}
\end{ex}

% ===== Câu 23 =====
% ID: L12C2B10-155-TLN
% Mức: VDC
\begin{ex}
% ID: L12C2B10-155-TLN
% Mức: VDC
Trong dạng “Bài toán tổng hợp về quá trình đẳng nhiệt”, Khí 28 lít ở $150\,\text{kPa}$ giãn/nén đẳng nhiệt đến $600\,\text{kPa}$. Tính thể tích theo lít.
\shortans{7}
\loigiai{
$V_2=150\cdot28/(600)=7$ lít.
}
\end{ex}

% ===== Câu 26 =====
% ID: L12C2B10-156-TL
% Mức: VD
\begin{bt}
% ID: L12C2B10-156-TL
% Mức: VD
Trong dạng “Khối khí chịu áp suất của cột chất lỏng”, Thể tích giảm từ 10 lít còn 5 lít ở nhiệt độ không đổi. Áp suất tăng bao nhiêu lần?
\loigiai{
$p_2/p_1=V_1/V_2=2$.
}
\end{bt}

% ===== Câu 27 =====
% ID: L12C2B10-157-TLN
% Mức: TH
\begin{ex}
% ID: L12C2B10-157-TLN
% Mức: TH
Trong dạng “Khối khí chịu áp suất của cột chất lỏng”, Khí 21 lít ở $140\,\text{kPa}$ được nén đẳng nhiệt còn 7 lít. Tính áp suất sau theo kPa.
\shortans{420}
\loigiai{
$p_2=140\cdot21/7=420$ kPa.
}
\end{ex}

% ===== Câu 30 =====
% ID: L12C2B10-158-TL
% Mức: VD
\begin{bt}
% ID: L12C2B10-158-TL
% Mức: VD
Trong dạng “Khối khí chịu áp suất của cột chất lỏng”, 110
\loigiai{
$p=770/7=110$ kPa.
}
\end{bt}

% ===== Câu 31 =====
% ID: L12C2B10-159-TLN
% Mức: TH
\begin{ex}
% ID: L12C2B10-159-TLN
% Mức: TH
Trong dạng “Khối khí chịu áp suất của cột chất lỏng”, Khí 3 lít ở $600\,\text{kPa}$ giãn đẳng nhiệt đến áp suất $150\,\text{kPa}$. Tính thể tích sau theo lít.
\shortans{12}
\loigiai{
$V_2=600\cdot3/150=12$ lít.
}
\end{ex}

% ===== Câu 34 =====
% ID: L12C2B10-160-TL
% Mức: VD
\begin{bt}
% ID: L12C2B10-160-TL
% Mức: VD
Trong dạng “Khối khí chịu áp suất của cột chất lỏng”, $\rho_2/\rho_1=p_2/p_1=4$.
\loigiai{
$\rho_2/\rho_1=p_2/p_1=4$.
}
\end{bt}

% ===== Câu 35 =====
% ID: L12C2B10-161-TLN
% Mức: VD
\begin{ex}
% ID: L12C2B10-161-TLN
% Mức: VD
Trong dạng “Khối khí chịu áp suất của cột chất lỏng”, Thể tích giảm từ 10 lít còn 5 lít ở nhiệt độ không đổi. Áp suất tăng bao nhiêu lần?
\shortans{2}
\loigiai{
$p_2/p_1=V_1/V_2=2$.
}
\end{ex}

% ===== Câu 39 =====
% ID: L12C2B10-162-TLN
% Mức: VD
\begin{ex}
% ID: L12C2B10-162-TLN
% Mức: VD
Trong dạng “Khối khí chịu áp suất của cột chất lỏng”, Một lượng khí có $pV=770$ kPa.lít. Khi V=7 lít, tính p theo kPa.
\shortans{110}
\loigiai{
$p=770/7=110$ kPa.
}
\end{ex}

% ===== Câu 43 =====
% ID: L12C2B10-163-TLN
% Mức: VD
\begin{ex}
% ID: L12C2B10-163-TLN
% Mức: VD
Trong dạng “Khối khí chịu áp suất của cột chất lỏng”, Áp suất đẳng nhiệt tăng từ $120\,\text{kPa}$ lên $480\,\text{kPa}$. Khối lượng riêng tăng bao nhiêu lần?
\shortans{4}
\loigiai{
$\rho_2/\rho_1=p_2/p_1=4$.
}
\end{ex}

% ===== Câu 47 =====
% ID: L12C2B10-164-TLN
% Mức: VDC
\begin{ex}
% ID: L12C2B10-164-TLN
% Mức: VDC
Trong dạng “Khối khí chịu áp suất của cột chất lỏng”, Khí 10 lít ở $130\,\text{kPa}$ giãn/nén đẳng nhiệt đến $260\,\text{kPa}$. Tính thể tích theo lít.
\shortans{5}
\loigiai{
$V_2=130\cdot10/(260)=5$ lít.
}
\end{ex}

% ===== Câu 50 =====
% ID: L12C2B10-165-TL
% Mức: VD
\begin{bt}
% ID: L12C2B10-165-TL
% Mức: VD
Trong dạng “Nhận biết quá trình đẳng nhiệt và định luật Boyle”, Khí 5 lít ở $120\,\text{kPa}$ được nén đẳng nhiệt còn 3 lít. Tính áp suất.
\loigiai{
$p_2=p_1V_1/V_2=120\cdot5/3=200$ kPa.
}
\end{bt}

% ===== Câu 51 =====
% ID: L12C2B10-166-TLN
% Mức: TH
\begin{ex}
% ID: L12C2B10-166-TLN
% Mức: TH
Trong dạng “Nhận biết quá trình đẳng nhiệt và định luật Boyle”, Khí có $p_1=100$ kPa, $V_1=6$ lít, nén đẳng nhiệt còn 4 lít. Tính $p_2$ theo kPa.
\shortans{150}
\loigiai{
$p_2=p_1V_1/V_2=100\cdot6/4=150$ kPa.
}
\end{ex}

% ===== Câu 55 =====
% ID: L12C2B10-167-TLN
% Mức: TH
\begin{ex}
% ID: L12C2B10-167-TLN
% Mức: TH
Trong dạng “Nhận biết quá trình đẳng nhiệt và định luật Boyle”, Khí ở 2 atm chiếm 3 lít, giãn đẳng nhiệt đến 1 atm. Tính thể tích theo lít.
\shortans{6}
\loigiai{
$V_2=p_1V_1/p_2=2\cdot3=6$ lít.
}
\end{ex}

% ===== Câu 56 =====
% ID: L12C2B10-168-TN
% Mức: TH
\begin{ex}
% ID: L12C2B10-168-TN
% Mức: TH
Trong dạng “Khối khí chịu áp suất của cột chất lỏng”, Khí có áp suất $120\,\text{kPa}$, thể tích 12 lít được nén đẳng nhiệt còn 3 lít. Áp suất sau là
\choice
{\True $480\,\text{kPa}$}
{$30\,\text{kPa}$}
{$120\,\text{kPa}$}
{$1920\,\text{kPa}$}
\loigiai{
Đáp án A. $p_2=p_1V_1/V_2=120\cdot12/3=480$ kPa.
}
\end{ex}

% ===== Câu 59 =====
% ID: L12C2B10-169-TLN
% Mức: VD
\begin{ex}
% ID: L12C2B10-169-TLN
% Mức: VD
Trong dạng “Nhận biết quá trình đẳng nhiệt và định luật Boyle”, Thể tích giảm từ 10 lít còn 2 lít ở nhiệt độ không đổi. Áp suất tăng bao nhiêu lần?
\shortans{5}
\loigiai{
$p_2/p_1=V_1/V_2=10/2=5$.
}
\end{ex}

% ===== Câu 60 =====
% ID: L12C2B10-170-TN
% Mức: TH
\begin{ex}
% ID: L12C2B10-170-TN
% Mức: TH
Trong dạng “Khối khí chịu áp suất của cột chất lỏng”, Khí có áp suất $125\,\text{kPa}$, thể tích 5 lít giãn đẳng nhiệt đến 25 lít. Áp suất sau là
\choice
{\True $25\,\text{kPa}$}
{$625\,\text{kPa}$}
{$125\,\text{kPa}$}
{$5\,\text{kPa}$}
\loigiai{
Đáp án A. $p_2=125\cdot5/25=25$ kPa.
}
\end{ex}

% ===== Câu 63 =====
% ID: L12C2B10-171-TLN
% Mức: VD
\begin{ex}
% ID: L12C2B10-171-TLN
% Mức: VD
Trong dạng “Nhận biết quá trình đẳng nhiệt và định luật Boyle”, Khí có $pV=800$ kPa.lít. Khi $V=5$ lít, tính p theo kPa.
\shortans{160}
\loigiai{
$p=800/5=160$ kPa.
}
\end{ex}

% ===== Câu 64 =====
% ID: L12C2B10-172-TN
% Mức: TH
\begin{ex}
% ID: L12C2B10-172-TN
% Mức: TH
Trong dạng “Khối khí chịu áp suất của cột chất lỏng”, Thể tích khí giảm 2 lần trong quá trình đẳng nhiệt thì áp suất
\choice
{\True tăng 2 lần}
{giảm 2 lần}
{không đổi}
{tăng 4 lần}
\loigiai{
Đáp án A. Do $p$ tỉ lệ nghịch với $V$.
}
\end{ex}

% ===== Câu 67 =====
% ID: L12C2B10-173-TLN
% Mức: VD
\begin{ex}
% ID: L12C2B10-173-TLN
% Mức: VD
Trong dạng “Nhận biết quá trình đẳng nhiệt và định luật Boyle”, Áp suất tăng từ $100\,\text{kPa}$ lên $250\,\text{kPa}$ đẳng nhiệt. Khối lượng riêng tăng bao nhiêu lần?
\shortans{2,5}
\loigiai{
$\rho_2/\rho_1=p_2/p_1=2,5$.
}
\end{ex}

% ===== Câu 71 =====
% ID: L12C2B10-174-TLN
% Mức: VDC
\begin{ex}
% ID: L12C2B10-174-TLN
% Mức: VDC
Trong dạng “Nhận biết quá trình đẳng nhiệt và định luật Boyle”, Khí 12 lít ở 1,5 atm được nén đẳng nhiệt đến 3 atm. Tính thể tích theo lít.
\shortans{6}
\loigiai{
$V_2=1,5\cdot12/3=6$ lít.
}
\end{ex}

% ===== Câu 74 =====
% ID: L12C2B10-175-TL
% Mức: VD
\begin{bt}
% ID: L12C2B10-175-TL
% Mức: VD
Trong dạng “Tính áp suất, thể tích trong quá trình đẳng nhiệt”, Khí 5 lít ở $120\,\text{kPa}$ được nén đẳng nhiệt còn 3 lít. Tính áp suất.
\loigiai{
$p_2=p_1V_1/V_2=120\cdot5/3=200$ kPa.
}
\end{bt}

% ===== Câu 75 =====
% ID: L12C2B10-176-TLN
% Mức: TH
\begin{ex}
% ID: L12C2B10-176-TLN
% Mức: TH
Trong dạng “Tính áp suất, thể tích trong quá trình đẳng nhiệt”, Khí có $p_1=100$ kPa, $V_1=6$ lít, nén đẳng nhiệt còn 4 lít. Tính $p_2$ theo kPa.
\shortans{150}
\loigiai{
$p_2=p_1V_1/V_2=100\cdot6/4=150$ kPa.
}
\end{ex}

% ===== Câu 79 =====
% ID: L12C2B10-177-TLN
% Mức: TH
\begin{ex}
% ID: L12C2B10-177-TLN
% Mức: TH
Trong dạng “Tính áp suất, thể tích trong quá trình đẳng nhiệt”, Khí ở 2 atm chiếm 3 lít, giãn đẳng nhiệt đến 1 atm. Tính thể tích theo lít.
\shortans{6}
\loigiai{
$V_2=p_1V_1/p_2=2\cdot3=6$ lít.
}
\end{ex}

% ===== Câu 83 =====
% ID: L12C2B10-178-TLN
% Mức: VD
\begin{ex}
% ID: L12C2B10-178-TLN
% Mức: VD
Trong dạng “Tính áp suất, thể tích trong quá trình đẳng nhiệt”, Thể tích giảm từ 10 lít còn 2 lít ở nhiệt độ không đổi. Áp suất tăng bao nhiêu lần?
\shortans{5}
\loigiai{
$p_2/p_1=V_1/V_2=10/2=5$.
}
\end{ex}

% ===== Câu 87 =====
% ID: L12C2B10-179-TLN
% Mức: VD
\begin{ex}
% ID: L12C2B10-179-TLN
% Mức: VD
Trong dạng “Tính áp suất, thể tích trong quá trình đẳng nhiệt”, Khí có $pV=800$ kPa.lít. Khi $V=5$ lít, tính p theo kPa.
\shortans{160}
\loigiai{
$p=800/5=160$ kPa.
}
\end{ex}

% ===== Câu 91 =====
% ID: L12C2B10-180-TLN
% Mức: VD
\begin{ex}
% ID: L12C2B10-180-TLN
% Mức: VD
Trong dạng “Tính áp suất, thể tích trong quá trình đẳng nhiệt”, Áp suất tăng từ $100\,\text{kPa}$ lên $250\,\text{kPa}$ đẳng nhiệt. Khối lượng riêng tăng bao nhiêu lần?
\shortans{2,5}
\loigiai{
$\rho_2/\rho_1=p_2/p_1=2,5$.
}
\end{ex}

% ===== Câu 98 =====
% ID: L12C2B10-181-TL
% Mức: VD
\begin{bt}
% ID: L12C2B10-181-TL
% Mức: VD
Trong dạng “Đồ thị đẳng nhiệt trên các hệ tọa độ”, Khí 5 lít ở $120\,\text{kPa}$ được nén đẳng nhiệt còn 3 lít. Tính áp suất.
\loigiai{
$p_2=p_1V_1/V_2=120\cdot5/3=200$ kPa.
}
\end{bt}

% ===== Câu 99 =====
% ID: L12C2B10-182-TLN
% Mức: TH
\begin{ex}
% ID: L12C2B10-182-TLN
% Mức: TH
Trong dạng “Đồ thị đẳng nhiệt trên các hệ tọa độ”, Khí có $p_1=100$ kPa, $V_1=6$ lít, nén đẳng nhiệt còn 4 lít. Tính $p_2$ theo kPa.
\shortans{150}
\loigiai{
$p_2=p_1V_1/V_2=100\cdot6/4=150$ kPa.
}
\end{ex}

% ===== Câu 103 =====
% ID: L12C2B10-183-TLN
% Mức: TH
\begin{ex}
% ID: L12C2B10-183-TLN
% Mức: TH
Trong dạng “Đồ thị đẳng nhiệt trên các hệ tọa độ”, Khí ở 2 atm chiếm 3 lít, giãn đẳng nhiệt đến 1 atm. Tính thể tích theo lít.
\shortans{6}
\loigiai{
$V_2=p_1V_1/p_2=2\cdot3=6$ lít.
}
\end{ex}

% ===== Câu 107 =====
% ID: L12C2B10-184-TLN
% Mức: VD
\begin{ex}
% ID: L12C2B10-184-TLN
% Mức: VD
Trong dạng “Đồ thị đẳng nhiệt trên các hệ tọa độ”, Thể tích giảm từ 10 lít còn 2 lít ở nhiệt độ không đổi. Áp suất tăng bao nhiêu lần?
\shortans{5}
\loigiai{
$p_2/p_1=V_1/V_2=10/2=5$.
}
\end{ex}

% ===== Câu 108 =====
% ID: L12C2B10-185-TN
% Mức: VD
\begin{ex}
% ID: L12C2B10-185-TN
% Mức: VD
Trong dạng “Nhận biết quá trình đẳng nhiệt và định luật Boyle”, Khí có $p_1=1$ atm, $V_1=4$ lít, nén còn 2 lít. $p_2$ bằng
\choice
{\True 2 atm}
{0,5 atm}
{4 atm}
{8 atm}
\loigiai{
Đáp án A. $p_2=p_1V_1/V_2=2$ atm.
}
\end{ex}

% ===== Câu 111 =====
% ID: L12C2B10-186-TLN
% Mức: VD
\begin{ex}
% ID: L12C2B10-186-TLN
% Mức: VD
Trong dạng “Đồ thị đẳng nhiệt trên các hệ tọa độ”, Khí có $pV=800$ kPa.lít. Khi $V=5$ lít, tính p theo kPa.
\shortans{160}
\loigiai{
$p=800/5=160$ kPa.
}
\end{ex}

% ===== Câu 115 =====
% ID: L12C2B10-187-TLN
% Mức: VD
\begin{ex}
% ID: L12C2B10-187-TLN
% Mức: VD
Trong dạng “Đồ thị đẳng nhiệt trên các hệ tọa độ”, Áp suất tăng từ $100\,\text{kPa}$ lên $250\,\text{kPa}$ đẳng nhiệt. Khối lượng riêng tăng bao nhiêu lần?
\shortans{2,5}
\loigiai{
$\rho_2/\rho_1=p_2/p_1=2,5$.
}
\end{ex}

% ===== Câu 119 =====
% ID: L12C2B10-188-TLN
% Mức: VDC
\begin{ex}
% ID: L12C2B10-188-TLN
% Mức: VDC
Trong dạng “Đồ thị đẳng nhiệt trên các hệ tọa độ”, Khí 12 lít ở 1,5 atm được nén đẳng nhiệt đến 3 atm. Tính thể tích theo lít.
\shortans{6}
\loigiai{
$V_2=1,5\cdot12/3=6$ lít.
}
\end{ex}

% ===== Câu 122 =====
% ID: L12C2B10-189-TL
% Mức: VD
\begin{bt}
% ID: L12C2B10-189-TL
% Mức: VD
Trong dạng “Ứng dụng định luật Boyle trong bơm, xi lanh và đời sống”, Thể tích giảm từ 18 lít còn 6 lít ở nhiệt độ không đổi. Áp suất tăng bao nhiêu lần?
\loigiai{
$p_2/p_1=V_1/V_2=3$.
}
\end{bt}

% ===== Câu 123 =====
% ID: L12C2B10-190-TLN
% Mức: TH
\begin{ex}
% ID: L12C2B10-190-TLN
% Mức: TH
\shortans{600}
\loigiai{
Áp dụng định luật Boyle-Mariotte cho quá trình nén đẳng nhiệt, ta có:
$p_1V_1 = p_2V_2$
Trong đó:
$p_1 = 150,	ext{kPa}$ (áp suất ban đầu)
$V_1 = 32,	ext{lít}$ (thể tích ban đầu)
$V_2 = 8,	ext{lít}$ (thể tích sau khi nén)
Ta cần tìm $p_2$ (áp suất sau khi nén).
Từ công thức trên, ta có:
$p_2 = \frac{p_1V_1}{V_2}$
Thay số vào, ta được:
$p_2 = \frac{150 \cdot 32}{8} = 150 \cdot 4 = 600,	ext{kPa}$
Vậy, áp suất của khối khí sau khi nén là $600,	ext{kPa}$.
}
\end{ex}

% ===== Câu 126 =====
% ID: L12C2B10-191-TL
% Mức: VD
\begin{bt}
% ID: L12C2B10-191-TL
% Mức: VD
Trong dạng “Ứng dụng định luật Boyle trong bơm, xi lanh và đời sống”, 120
\loigiai{
$p=960/8=120$ kPa.
}
\end{bt}

% ===== Câu 127 =====
% ID: L12C2B10-192-TLN
% Mức: TH
\begin{ex}
% ID: L12C2B10-192-TLN
% Mức: TH
Trong dạng “Ứng dụng định luật Boyle trong bơm, xi lanh và đời sống”, Khí 4 lít ở $200\,\text{kPa}$ giãn đẳng nhiệt đến áp suất $100\,\text{kPa}$. Tính thể tích sau theo lít.
\shortans{8}
\loigiai{
$V_2=200\cdot4/100=8$ lít.
}
\end{ex}

% ===== Câu 130 =====
% ID: L12C2B10-193-TL
% Mức: VD
\begin{bt}
% ID: L12C2B10-193-TL
% Mức: VD
Trong dạng “Ứng dụng định luật Boyle trong bơm, xi lanh và đời sống”, $\rho_2/\rho_1=p_2/p_1=2$.
\loigiai{
$\rho_2/\rho_1=p_2/p_1=2$.
}
\end{bt}

% ===== Câu 131 =====
% ID: L12C2B10-194-TLN
% Mức: VD
\begin{ex}
% ID: L12C2B10-194-TLN
% Mức: VD
Trong dạng “Ứng dụng định luật Boyle trong bơm, xi lanh và đời sống”, Thể tích giảm từ 18 lít còn 6 lít ở nhiệt độ không đổi. Áp suất tăng bao nhiêu lần?
\shortans{3}
\loigiai{
$p_2/p_1=V_1/V_2=3$.
}
\end{ex}

% ===== Câu 134 =====
% ID: L12C2B10-195-TLN
% Mức: VD
\begin{ex}
% ID: L12C2B10-195-TLN
% Mức: VD
Trong dạng “Ứng dụng định luật Boyle trong bơm, xi lanh và đời sống”, Một lượng khí có $pV=960$ kPa.lít. Khi V=8 lít, tính p theo kPa.
\shortans{120}
\loigiai{
$p=960/8=120$ kPa.
}
\end{ex}

% ===== Câu 136 =====
% ID: L12C2B10-196-TLN
% Mức: VD
\begin{ex}
% ID: L12C2B10-196-TLN
% Mức: VD
Trong dạng “Ứng dụng định luật Boyle trong bơm, xi lanh và đời sống”, Áp suất đẳng nhiệt tăng từ $130\,\text{kPa}$ lên $260\,\text{kPa}$. Khối lượng riêng tăng bao nhiêu lần?
\shortans{2}
\loigiai{
$\rho_2/\rho_1=p_2/p_1=2$.
}
\end{ex}

% ===== Câu 138 =====
% ID: L12C2B10-197-TLN
% Mức: VDC
\begin{ex}
% ID: L12C2B10-197-TLN
% Mức: VDC
Trong dạng “Ứng dụng định luật Boyle trong bơm, xi lanh và đời sống”, Khí 18 lít ở $140\,\text{kPa}$ giãn/nén đẳng nhiệt đến $420\,\text{kPa}$. Tính thể tích theo lít.
\shortans{6}
\loigiai{
$V_2=140\cdot18/(420)=6$ lít.
}
\end{ex}

% ===== Câu 140 =====
% ID: L12C2B10-198-TN
% Mức: VD
\begin{ex}
% ID: L12C2B10-198-TN
% Mức: VD
Trong dạng “Tính áp suất, thể tích trong quá trình đẳng nhiệt”, Khí có $p_1=1$ atm, $V_1=4$ lít, nén còn 2 lít. $p_2$ bằng
\choice
{\True 2 atm}
{0,5 atm}
{4 atm}
{8 atm}
\loigiai{
Đáp án A. $p_2=p_1V_1/V_2=2$ atm.
}
\end{ex}

% ===== Câu 150 =====
% ID: L12C2B10-199-TN
% Mức: VD
\begin{ex}
% ID: L12C2B10-199-TN
% Mức: VD
Trong dạng “Đồ thị đẳng nhiệt trên các hệ tọa độ”, Khí có $p_1=1$ atm, $V_1=4$ lít, nén còn 2 lít. $p_2$ bằng
\choice
{\True 2 atm}
{0,5 atm}
{4 atm}
{8 atm}
\loigiai{
Đáp án A. $p_2=p_1V_1/V_2=2$ atm.
}
\end{ex}

% ===== Câu 157 =====
% ID: L12C2B10-200-TN
% Mức: TH
\begin{ex}
% ID: L12C2B10-200-TN
% Mức: TH
Trong dạng “Ứng dụng định luật Boyle trong bơm, xi lanh và đời sống”, Khí có áp suất $125\,\text{kPa}$, thể tích 20 lít được nén đẳng nhiệt còn 4 lít. Áp suất sau là
\choice
{\True $625\,\text{kPa}$}
{$25\,\text{kPa}$}
{$125\,\text{kPa}$}
{$3125\,\text{kPa}$}
\loigiai{
Đáp án A. $p_2=p_1V_1/V_2=125\cdot20/4=625$ kPa.
}
\end{ex}

% ===== Câu 158 =====
% ID: L12C2B10-201-TN
% Mức: TH
\begin{ex}
% ID: L12C2B10-201-TN
% Mức: TH
\choice
{\True $65\,\text{kPa}$}
{$260\,\text{kPa}$}
{$130\,\text{kPa}$}
{$32,5\,\text{kPa}$}
\loigiai{
Áp dụng định luật Boyle-Mariotte cho quá trình đẳng nhiệt, ta có $p_1V_1 = p_2V_2$.
Với $p_1 = 130\,\text{kPa}, V_1 = 6\,\text{lít}$ và $V_2 = 12\,\text{lít}$.
Ta có $p_2 = \frac{p_1V_1}{V_2} = \frac{130\cdot6}{12} = 65\,\text{kPa}$.
Vậy, áp suất sau khi giãn nở là $65\,\text{kPa}$.
}
\end{ex}
